\section{Query Engine}
\label{sec:query-engine}

The Query Engine translates structured MCP tool requests into SQL queries
against the symbol graph. This section describes the major query patterns.
Motivation for tool-based code navigation is grounded in empirical studies of
how developers investigate large codebases~\cite{robillard2004codenav}.

\subsection{Symbol search}

\texttt{symbol\_search} queries the \texttt{nodes\_fts} FTS5 virtual table
(\S\ref{sec:data-model}) with a prefix \texttt{MATCH} on the \texttt{name}
column:

\begin{lstlisting}[language=sql]
SELECT n.id, n.name, n.qualname, n.kind, f.path, n.span
FROM   nodes n
JOIN   files f ON f.id = n.file_id
JOIN   nodes_fts nf ON nf.rowid = n.id
WHERE  nf.name MATCH :query || '*'
  AND  n.node_type = 'symbol'
  AND  (:root_id IS NULL OR f.root_id = :root_id)
ORDER  BY rank
LIMIT  :limit;
\end{lstlisting}

A \texttt{kind='function'} filter expands to \texttt{kind IN ('function',
'method')} to achieve polyglot parity: in Python a method has kind
\texttt{method} while in C a standalone function has kind \texttt{function},
but both are callable symbols.

\subsection{Call graph traversal}
\label{sec:qe-callgraph}

Call graph construction at query time is based on edges extracted by the
indexer~\cite{ryder1979callgraph,grove2001callgraph}. Because codetopo
does not perform type resolution, these edges are AST-derived approximations
--- hence the ``approx'' suffix on the tool names.

\subsubsection{callers\_approx}

\begin{lstlisting}[language=sql]
SELECT n.id, n.name, n.kind, f.path, n.span
FROM   edges e
JOIN   nodes n ON n.id = e.src_id
JOIN   files f ON f.id = n.file_id
WHERE  e.dst_id = :node_id
  AND  e.kind   = 'calls';
\end{lstlisting}

\subsubsection{callees\_approx}

\begin{lstlisting}[language=sql]
SELECT e.dst_id, e.dst_name, n.kind, f.path, n.span
FROM   edges e
LEFT   JOIN nodes n ON n.id = e.dst_id
LEFT   JOIN files f ON f.id = n.file_id
WHERE  e.src_id = :node_id
  AND  e.kind   = 'calls';
\end{lstlisting}

When \texttt{dst\_id} is \texttt{NULL}, the callee is returned by name only
(\texttt{dst\_name}): an unresolved reference retained from the extractor.

\subsection{Impact analysis}

\texttt{impact\_of} computes the transitive closure of the call graph starting
from a seed node using a BFS/DFS over \texttt{calls} edges. A recursive
CTE expresses this efficiently in SQLite:

\begin{lstlisting}[language=sql]
WITH RECURSIVE impact(id, depth) AS (
  SELECT :seed_id, 0
  UNION ALL
  SELECT e.dst_id, i.depth + 1
  FROM   edges e
  JOIN   impact i ON e.src_id = i.id
  WHERE  e.kind = 'calls' AND i.depth < :max_depth
)
SELECT DISTINCT n.id, n.name, n.kind, f.path, i.depth
FROM   impact i
JOIN   nodes n ON n.id = i.id
JOIN   files f ON f.id = n.file_id
ORDER  BY i.depth, n.name;
\end{lstlisting}

\texttt{shortest\_path} uses a similar CTE with early termination when the
target node is reached.

\subsection{Path resolution}

\texttt{resolve\_db\_path()} converts a filesystem path into a database
\texttt{file\_id}. It tries three strategies:

\begin{enumerate}
  \item \textbf{Absolute match}: \texttt{roots.path || '/' || files.path = :abs}
  \item \textbf{Exact relative match}: \texttt{files.path = :rel}
  \item \textbf{Suffix LIKE match}: \texttt{files.path LIKE '\%' || :suffix}
\end{enumerate}

Strategy 3 handles workspace roots with absolute paths that differ between
machines, and editor integrations that pass workspace-relative paths.

\subsection{Content full-text search}

\texttt{code\_search} uses \texttt{content\_fts}, the contentless trigram FTS5
table populated per source line~\cite{cox2012trigram}:

\begin{lstlisting}[language=sql]
SELECT f.path, r.path AS root_path,
       cft.line_no,
       snippet(content_fts, 0, '', '', '...', 10) AS excerpt
FROM   content_fts
JOIN   content_fts_tracker cft USING (file_id)
JOIN   files f ON f.id = cft.file_id
JOIN   roots r ON r.id = f.root_id
WHERE  content_fts MATCH :query
  AND  (:root_id IS NULL OR f.root_id = :root_id)
ORDER  BY rank
LIMIT  :limit;
\end{lstlisting}

Because the table is contentless, \texttt{content\_fts\_tracker} is joined to
retrieve the file ID and line number, which are stored as unindexed columns in
the FTS row.
